\section{Evaluasi Model}

Evaluasi performa model merupakan tahapan penting dalam pembelajaran mesin (\textit{machine learning}) untuk menilai seberapa baik model mampu melakukan prediksi terhadap data yang belum pernah dilihat sebelumnya. Menurut Tan, Steinbach, dan Kumar (2019) dalam buku \textit{Introduction to Data Mining}, beberapa metrik umum yang digunakan dalam evaluasi model klasifikasi mencakup akurasi, presisi, recall, dan F1-score. Masing-masing metrik memberikan perspektif berbeda terhadap performa model, sehingga pemilihan metrik yang tepat bergantung pada konteks dan karakteristik data.
    
    \begin{itemize}
        \item {Akurasi} (\textit{accuracy}) adalah rasio jumlah prediksi yang benar terhadap total prediksi.
        \[
        \text{Akurasi} = \frac{TP + TN}{TP + TN + FP + FN}
        \]
        dengan:
        \begin{itemize}
            \item $TP$ (True Positive): jumlah kasus positif yang berhasil diprediksi dengan benar,
            \item $TN$ (True Negative): jumlah kasus negatif yang berhasil diprediksi dengan benar,
            \item $FP$ (False Positive): jumlah kasus negatif yang salah diklasifikasikan sebagai positif,
            \item $FN$ (False Negative): jumlah kasus positif yang salah diklasifikasikan sebagai negatif.
        \end{itemize}
        Meskipun akurasi memberikan gambaran umum performa model, metrik ini dapat menyesatkan pada dataset yang tidak seimbang. Misalnya, jika hanya 5\% dari total data merupakan penyewa mencurigakan, model yang selalu memprediksi ``normal'' tetap bisa memiliki akurasi tinggi.
    
        \item {Presisi} (\textit{precision}) mengukur ketepatan prediksi positif, yaitu seberapa besar proporsi prediksi positif yang benar-benar positif.
        \[
        \text{Presisi} = \frac{TP}{TP + FP}
        \]
        Presisi tinggi menunjukkan bahwa model jarang menghasilkan alarm palsu. Dalam konteks deteksi penyewa mencurigakan, presisi penting karena kesalahan dalam menandai penyewa normal sebagai mencurigakan dapat berdampak sosial dan hukum.
    
        \item {Recall} (atau \textit{sensitivity}) mengukur seberapa baik model dalam menemukan seluruh kasus positif.
        \[
        \text{Recall} = \frac{TP}{TP + FN}
        \]
        Recall tinggi menunjukkan bahwa model berhasil menangkap hampir semua penyewa yang benar-benar mencurigakan. Namun, jika hanya fokus pada recall, model bisa menghasilkan banyak alarm palsu sehingga presisinya rendah.
    
        \item {F1-score} adalah rata-rata harmonis dari presisi dan recall.
        \[
        \text{F1-score} = 2 \cdot \frac{\text{Presisi} \cdot \text{Recall}}{\text{Presisi} + \text{Recall}}
        \]
        Nilai F1-score berada antara 0 hingga 1. Nilai mendekati 1 menunjukkan bahwa model memiliki keseimbangan yang baik antara presisi dan recall, sehingga cocok digunakan ketika keduanya sama-sama penting.
    \end{itemize}
    
    

    Untuk menggambarkan kelemahan metrik akurasi pada dataset yang tidak seimbang, literatur standar pembelajaran mesin sering menggunakan skenario deteksi anomali, seperti penipuan kartu kredit atau deteksi penyakit langka (Provost \& Fawcett, 2013). Sebagai ilustrasi yang diadaptasi dari studi kasus evaluasi model klasifikasi, misalkan terdapat 10.000 data transaksi (dengan analogi data penyewa), di mana hanya 100 di antaranya adalah kasus kecurangan nyata (mencurigakan), sedangkan 9.900 sisanya adalah transaksi normal. Sebuah model deteksi mencurigai 150 transaksi, dan dari jumlah tersebut sebesar 90 adalah benar-benar kecurangan, maka:
    
    \begin{itemize}
        \item $TP = 90$ (terdeteksi mencurigakan dan benar kecurangan)
        \item $FP = 60$ (terdeteksi mencurigakan, tetapi sebenarnya normal)
        \item $FN = 10$ (tidak terdeteksi, padahal kecurangan)
        \item $TN = 9840$ (tidak terdeteksi dan memang normal)
    \end{itemize}
            
    Berikut perhitungannya.
            
    \[
    \text{Akurasi} = \frac{TP + TN}{10000} = \frac{90 + 9840}{10000} = 0{,}993 \quad (99{,}3\%)
    \]
    \[
    \text{Presisi} = \frac{90}{90 + 60} = \frac{90}{150} = 0{,}60 \quad (60\%)
    \]
    \[
    \text{Recall} = \frac{90}{90 + 10} = \frac{90}{100} = 0{,}90 \quad (90\%)
    \]
    \[
    \text{F1-score} = 2 \cdot \frac{0{,}60 \cdot 0{,}90}{0{,}60 + 0{,}90} = \frac{1{,}08}{1{,}50} = 0{,}72 \quad (72\%)
    \]
            
    Hasil evaluasi ini menunjukkan fenomena yang dikenal sebagai Paradoks Akurasi (\textit{Accuracy Paradox}). Meskipun model seolah-olah berkinerja sangat luar biasa dengan akurasi 99,3\%, nilai presisinya hanya 60\%. Hal ini berarti 40\% dari total peringatan yang dikeluarkan oleh sistem adalah alarm palsu (\textit{false positive}). Dalam konteks penyewaan properti, jika pengelola hanya melihat nilai akurasi, mereka mungkin menganggap sistem sudah sempurna. Namun, presisi 60\% menunjukkan risiko tinggi di mana banyak penyewa normal yang baik-baik saja akan dituduh atau ditandai sebagai penyewa mencurigakan. Kesalahan identifikasi ini dapat memicu konflik dengan penyewa, pelanggaran privasi, atau hilangnya pendapatan penyewaan. 
    
    Kelemahan metrik ini sejalan dengan temuan Fawcett (2006) dalam \textit{An Introduction to ROC Analysis}, yang menyatakan bahwa akurasi cenderung tidak dapat diandalkan ketika jumlah kasus positif sangat sedikit. Lebih lanjut, menurut Saito dan Rehmsmeier (2015) dalam jurnal \textit{PLOS ONE}, metrik presisi dan \textit{recall} (termasuk \textit{F1-score}) terbukti lebih informatif dan memberikan gambaran performa yang jauh lebih jujur dibandingkan akurasi pada dataset yang tidak seimbang. Maka dari itu, evaluasi performa model dalam penelitian ini akan lebih menitikberatkan pada metrik presisi, guna meminimalkan tingkat \textit{false positive} dan menjaga keandalan sistem dalam mendeteksi penyewa mencurigakan.





    
    